\documentclass[../main.tex]{subfiles}

\begin{document}

\chapter{一般测度空间: 性质与构造}\label{chapter:ch17-general-measure}

本章的第一个目标是抽象出没有任何拓扑的实直线上的 Lebesgue 测度的最重要性质. 通过给出 Lebesgue 测度满足的某些公理, 将我们的理论建立在这些公理的基础上. 因此我们的理论适用于每个满足这些给定公理的系统.

为了证明实直线上的 Lebesgue 测度是 $\sigma$ 代数上的可数可加的集函数, 我们仅用了最为初等的集论概念. 通过赋予每个有界区间长度, 我们定义了原始的集函数. 将此集函数延拓为对实数的每个子集有定义的外测度集函数后, 我们区分出可测集族. 我们证明可测集是 $\sigma$ 代数. 外测度在其上的限制是一个测度. 我们称此为 Lebesgue 测度的 Carath\'eodory 构造法. 本章的第二个目标是证明对一般的抽象集合 $X$, Carath\'eodory 构造法是可行的. 事实上, 我们证明任何定义在 $X$ 的子集族 $\mathcal{S}$ 上的非负集函数 $\mu$ 诱导出一个外测度 $\mu^*$. 关于该外测度我们能够确定可测集的 $\sigma$ 代数 $\mathcal{M}$. $\mu^*$ 在 $\mathcal{M}$ 上的限制是一个测度, 我们称它为由 $\mu$ 诱导的 Carath\'eodory 测度. 我们以 Carath\'eodory-Hahn 定理的证明结束本章, 该定理告诉我们在非常一般的条件下由集函数 $\mu$ 诱导的 Carath\'eodory 测度是 $\mu$ 的一个延拓.

\subfile{Chapter-17/section-01.tex}

\subfile{Chapter-17/section-02.tex}

\subfile{Chapter-17/section-03.tex}

\subfile{Chapter-17/section-04.tex}

\subfile{Chapter-17/section-05.tex}

\end{document}